% Generated draft: requires mathematical and source review.
\documentclass[11pt]{article}
\usepackage[margin=25mm]{geometry}
\usepackage{fontspec}
\setmainfont{Latin Modern Roman}
\usepackage{amsmath,amssymb,mathtools,unicode-math}
\setmathfont{Latin Modern Math}
\usepackage{enumitem,needspace,xurl,hyperref}
\setlength{\parindent}{0pt}
\setlength{\parskip}{5pt}
\setlength{\emergencystretch}{3em}
\newcommand{\N}{\mathbb{N}}
\newcommand{\Z}{\mathbb{Z}}
\newcommand{\Q}{\mathbb{Q}}
\newcommand{\R}{\mathbb{R}}
\newcommand{\C}{\mathbb{C}}
\newcommand{\F}{\mathbb{F}}
\newcommand{\A}{\mathbb{A}}
\newcommand{\PP}{\mathbb{P}}
\newcommand{\E}{\mathbb{E}}
\newcommand{\eps}{\varepsilon}
\newcommand{\dd}{\,\mathrm{d}}
\begin{document}
\section*{UnsolvedMath ID: 30004966}
\textbf{Source problem code:} catalogue entry
\subsection{Definitions and mathematical statement}
For a graph \(H\) and a graph \(G\), let
\[
N_H(G)=\#\{F:F\text{ is a subgraph of }G\text{ and }F\cong H\}.
\]
This follows the supplied clean statement's wording of subgraphs isomorphic to \(H\); the counting convention and permitted graph types require confirmation from the source. For a graph class \(\mathcal G\) and a nonnegative integer \(n\), write
\[
\mathcal G_n=\{G\in\mathcal G:|V(G)|=n\}.
\]
Whenever \(\mathcal G_n\) is nonempty and the indicated maximum exists, the requested extremal quantity is
\[
\operatorname{ex}(H,\mathcal G,n)
=\max_{G\in\mathcal G_n}N_H(G).
\]
\textbf{Requested determination.} Determine \(\operatorname{ex}(H,\mathcal G,n)\) for the intended graph \(H\), graph class \(\mathcal G\), and admissible vertex counts \(n\).

\textbf{Unresolved qualifications.} The supplied record specifies neither \(H\) nor \(\mathcal G\), and supplies no definition of the robust sparsity mentioned in its contextual talk title. It also leaves graph types, counting conventions, the range of \(n\), and the treatment of \(\mathcal G_n=\varnothing\) unspecified. Consequently this is a conditional description of the requested extremal quantity, not a fully specified extremal problem. No value is assigned when the defining maximum is unavailable.
\subsection{Short English statement}
The question asks for the largest number of copies of a graph \(H\) that can occur in an \(n\)-vertex graph belonging to a class \(\mathcal G\). The supplied record does not identify \(H\) or \(\mathcal G\), so the intended concrete problem remains unresolved. Its open status is a supplied curation label.
\subsection{Sources}
\begin{itemize}
\item Local UnsolvedMath registry, permanent ID \(30004966\), title \emph{Extremal Subgraph Counts in Sparse Graph Classes}, source problem code \texttt{catalogue entry}; collection: Oberwolfach Reports Open Problems.
\item Supplied source citation: \emph{Graph Theory (2023)}, Oberwolfach Reports, DOI \texttt{10.4171/owr/2022/1}, pp.~5--78. Supplied source URL: \url{https://doi.org/10.4171/owr/2022/1}. The supplied literature reference labels this same link \emph{Oberwolfach Report 1/2022}.
\item The original extracted statement appends the contextual title \emph{Homomorphism counts in robustly sparse graphs} and the name Chun-Hung Liu. The record does not explicitly identify a problem proposer; its proposed-year field is \(2022\).
\item Literature assessment dated \texttt{2026-08-22}: supplied dated curation, not independently verified here. It labels the problem open because the extracted record does not define \(\mathcal G\) or \(H\), preventing verification of a determinate extremal claim. Its listed source is \emph{Oberwolfach Report 1/2022}, \url{https://doi.org/10.4171/owr/2022/1}. The background requests independent verification of time-sensitive status and cited claims.
\item The record labels the statement \texttt{corrected\_verified} and reports verification against the official Oberwolfach report. This is supplied verification metadata, not an independent verification performed for this draft.
\end{itemize}
\end{document}
